% ===== PRÉAMBULE CANONIQUE — à coller VERBATIM en tête de CHAQUE .tex =====
\documentclass[11pt,a4paper]{article}
\usepackage[utf8]{inputenc}\usepackage[T1]{fontenc}\usepackage{lmodern}
\usepackage[french]{babel}
\usepackage{amsmath,amssymb,mathtools}\usepackage{icomma}
\usepackage{siunitx}\sisetup{locale=FR}  % \qty{}{}, \si{}, \ang{}
\usepackage{xcolor}\usepackage{colortbl}
\usepackage{tikz}\usepackage{pgfplots}\pgfplotsset{compat=1.18}
\usepackage[siunitx,europeanresistors]{circuitikz} % circuits (élec.)
\usepackage[most]{tcolorbox}\usepackage{enumitem}\usepackage{array,booktabs}
\usepackage[a4paper,margin=2.2cm]{geometry}
\usetikzlibrary{arrows.meta,calc,angles,quotes,positioning,patterns,%
  backgrounds,shapes.geometric,decorations.pathmorphing,decorations.markings,intersections}
\definecolor{primary}{RGB}{222,49,99}\definecolor{primarydark}{RGB}{150,22,55}
\definecolor{defcol}{RGB}{0,114,178}\definecolor{methcol}{RGB}{52,92,148}
\definecolor{excol}{RGB}{0,135,90}\definecolor{attncol}{RGB}{200,40,55}
\tcbset{boxbase/.style={enhanced,breakable,boxrule=0.9pt,arc=2.5pt,
  left=3mm,right=3mm,top=2.3mm,bottom=2.3mm,fonttitle=\bfseries,coltitle=white}}
% --- boîtes TUTORIEL (noms EXACTS, ne pas renommer) ---
\newtcolorbox{cle}[1][]{boxbase,colback=defcol!6,colframe=defcol,
  title={\textsc{À retenir}\ifx\relax#1\relax\relax\else\textmd{ : \itshape #1}\fi}}
\newtcolorbox{methode}[1][]{boxbase,colback=methcol!7,colframe=methcol,sharp corners,
  title={\textsc{Méthode}\ifx\relax#1\relax\relax\else\textmd{ --- \itshape #1}\fi}}
\newtcolorbox{exemplebox}[1][]{boxbase,colback=excol!5,colframe=excol,
  title={\textsc{Exemple}\ifx\relax#1\relax\relax\else\textmd{ : \itshape #1}\fi}}
\newtcolorbox{piege}[1][]{boxbase,colback=attncol!6,colframe=attncol,
  title={\textsc{Piège}\ifx\relax#1\relax\relax\else\textmd{ : \itshape #1}\fi}}
\newtcolorbox{remarque}[1][]{boxbase,colback=black!4,colframe=black!45,
  title={\textsc{Remarque}\ifx\relax#1\relax\relax\else\textmd{ : \itshape #1}\fi}}
% --- boîtes EXERCICE / CORRIGÉ (noms EXACTS, auto-comptées) ---
\newtcolorbox[auto counter]{exo}[1][]{boxbase,colback=primary!4,colframe=primary,
  title={\textsc{Exercice}~\thetcbcounter\ifx\relax#1\relax\relax\else\textmd{ --- \itshape #1}\fi}}
\newtcolorbox[auto counter]{corrige}[1][]{boxbase,colback=excol!4,colframe=excol,
  title={\textsc{Corrigé}~\thetcbcounter\ifx\relax#1\relax\relax\else\textmd{ --- \itshape #1}\fi}}
\newcommand{\vv}[1]{\vec{#1}}\newcommand{\ds}{\displaystyle}
\newcommand{\R}{\mathbb{R}}\newcommand{\N}{\mathbb{N}}\newcommand{\Z}{\mathbb{Z}}
% (un \usepackage{fancyhdr}+en-tête est possible ; il est ignoré côté web)
% =========================================================================
\begin{document}

\begin{center}
{\large\bfseries\color{primarydark} Thème 2 — Niveau Facile}\\[-2pt]
{\color{primary}\rule{5cm}{1pt}}
\end{center}

\begin{exo}[calculer la pulsation $\omega$]
On rappelle $\omega=2\pi f=2\pi/T$.
\begin{enumerate}[label=\alph*)]
  \item Une onde a pour fréquence $f=\SI{10}{\hertz}$. Calcule sa pulsation $\omega$
        (valeur exacte en fonction de $\pi$, puis valeur approchée en \si{\radian\per\second}).
  \item Une autre onde a pour période $T=\SI{0.5}{\second}$. Calcule sa pulsation $\omega$.
\end{enumerate}
\end{exo}

\begin{exo}[calculer le nombre d'onde $k$]
On rappelle $k=2\pi/\lambda$.
\begin{enumerate}[label=\alph*)]
  \item Une onde a pour longueur d'onde $\lambda=\SI{0.5}{\metre}$. Calcule $k$.
  \item Une autre a $\lambda=\SI{0.2}{\metre}$. Calcule $k$.
\end{enumerate}
\end{exo}

\begin{exo}[lire les paramètres dans l'équation]
Une onde est décrite (unités SI) par
\[ y(x,t)=0{,}05\,\sin\!\left(10\pi\,t - 4\pi\,x\right). \]
\begin{enumerate}[label=\alph*)]
  \item Donne l'amplitude $A$.
  \item Donne la pulsation $\omega$ et le nombre d'onde $k$.
  \item L'onde se propage-t-elle vers les $x$ croissants ou décroissants ?
\end{enumerate}
\end{exo}

\begin{exo}[vitesse et accélération maximales d'un point]
Un point d'une corde oscille avec l'amplitude $A=\SI{0.02}{\metre}$ et la pulsation
$\omega=\SI{50}{\radian\per\second}$.
\begin{enumerate}[label=\alph*)]
  \item Calcule sa vitesse maximale $v_{\max}=A\omega$.
  \item Calcule son accélération maximale $a_{\max}=A\omega^{2}$.
\end{enumerate}
\end{exo}

\begin{exo}[concordance ou opposition de phase ?]
La photo ci-dessous montre une onde de longueur d'onde $\lambda=\SI{4}{\metre}$. Trois points
$M$, $N$, $P$ sont repérés.
\begin{center}
\begin{tikzpicture}
\begin{axis}[
  width=12cm,height=4.8cm,
  xmin=0,xmax=10,ymin=-2.4,ymax=2.4,
  xtick={0,1,2,3,4,5,6,7,8,9,10},ytick=\empty,
  grid=both,grid style={black!12},axis lines=middle,
  xlabel={$x$ (m)},
  xlabel style={at={(axis description cs:0.5,-0.02)},anchor=north},
  tick label style={font=\footnotesize},samples=200,clip=false]
  \addplot[defcol,line width=1.3pt,domain=0:10]{2*sin(deg(2*pi*x/4))};
  \addplot[only marks,mark=*,primary,mark size=2.4pt] coordinates {(1,2)(3,-2)(5,2)};
  \node[above,primary,font=\bfseries] at (axis cs:1,2){$M$};
  \node[below,primary,font=\bfseries] at (axis cs:3,-2){$N$};
  \node[above,primary,font=\bfseries] at (axis cs:5,2){$P$};
\end{axis}
\end{tikzpicture}
\end{center}
Pour chaque paire, dis si les points vibrent en \textbf{concordance} ou en \textbf{opposition}
de phase (calcule d'abord la distance qui les sépare).
\begin{enumerate}[label=\alph*)]
  \item $M$ et $P$.
  \item $M$ et $N$.
\end{enumerate}
\end{exo}

\end{document}
